\chapter{Сборка полного примитивного квантово-марковского генератора}
\label{ch:v8-full-primitive-markov-generator-assembly-gate}

\section{Один процесс из ранее выведенных частей}

Предыдущие главы по отдельности получили связывающую диффузию, калибровочный селектор
и межсекторный Краус-мост. Теперь они собираются в один генератор
\begin{equation}
 \mathcal L_{\rm full}
 =\mathcal L_{\rm link}
 +\mathcal L_{SU(3)}
 +\mathcal L_{SU(2)}
 +\mathcal L_{U(1)}
 +\mathcal L_{QLYR}
 +\mathcal L_{XLdR}.
 \label{eq:v8-full-primitive-generator}
\end{equation}
Все шесть слагаемых являются симметричными формами Дирихле на уже
выведенной концевой алгебре
\begin{equation}
 \mathcal A_{\rm obs}=M_{11}(\mathbb C)\oplus M_{10}(\mathbb C).
\end{equation}

\section{Полная положительность и примитивность}

Генератор \eqref{eq:v8-full-primitive-generator} имеет явное представление
Линдблада с 25 самосопряжёнными операторами скачка: один связывающий оператор,
двенадцать калибровочных генераторов и двенадцать межсекторных стрелок операторов. Поэтому
$\exp(t\mathcal L_{\rm full})$ является вполне положительной унитальной
полугруппой. Машинная сборка полного $221\times221$ ограничения проверяет
сохранение концевой алгебры, следа, единицы и полугрупповое тождество.

Главный результат сборки:
\begin{equation}
 \ker\mathcal L_{\rm full}=\mathbb CI_{21}.
 \label{eq:v8-full-generator-scalar-kernel}
\end{equation}
Без межсекторной части ядро имеет размерность два. Добавление любой одной полной
цветовой межсекторные семьи уже соединяет два центральных сектора. Значит,
совместный процесс примитивен и имеет единственное верное стационарное
состояние $I_{21}/21$.

В незавешенном представителе имеется строго положительная спектральная
щель. Случайный скан 48 наборов независимых положительных весов всех шести
слагаемых в диапазоне $10^{-3}\ldots10^3$ всегда оставляет ядро
одномерным и спектр неположительным. Получен не один случайный процесс, а
качественный универсальный класс примитивных полугрупп.

\section{Что подробный баланс уже может и не может выбрать}

Каждое слагаемое самосопряжено относительно скалярного произведения
Гильберта--Шмидта. Поэтому любое положительно взвешенное семейство
\begin{equation}
 \mathcal L_{\boldsymbol\kappa}
 =\sum_r\kappa_r\mathcal L_r,
 \qquad \kappa_r>0,
 \label{eq:v8-positive-rate-family}
\end{equation}
удовлетворяет подробный баланс относительно нормированного следа. Следовая
обратимость гарантирует правильное равновесие, но не выбирает шесть
относительных скоростей. Спектральная щель заметно меняется при их
изменении, хотя примитивность сохраняется.

\begin{theorem}[Примитивный марковский универсальный класс]
Уже выведенные данные связывания, калибровки и межсекторных стрелок порождают на одной
концевой алгебре вполне положительную, унитальную, сохраняющую след
квантово-марковскую полугруппу с неподвижной алгеброй $\mathbb CI_{21}$.
Примитивность устойчива при любых строго положительных относительных весах.
Следовательно, качественная сшивка наблюдаемых завершена; незамкнутой
остаётся метрика скоростей.
\end{theorem}

\begin{nogocriterion}[Запрет объявлять детальный баланс по следу селектором]
Нельзя использовать следовую обратимость для выбора коэффициентов
$\kappa_r$: каждое слагаемое удовлетворяет ей отдельно. Для количественного
замыкания требуется нетривиальное KMS-состояние, микроскопическая
ковариация либо пространственно-временная нормировка.
\end{nogocriterion}

\begin{protocol}
\begin{enumerate}
 \item единая алгебра наблюдаемых: \textbf{$M_{11}\oplus M_{10}$};
 \item самосопряжённых операторов скачка: \textbf{25};
 \item полная положительность: \textbf{сертифицирована формой Линдблада};
 \item сохранение единицы и следа: \textbf{пройдено};
 \item неподвижная алгебра: \textbf{$\mathbb CI_{21}$};
 \item примитивность: \textbf{получена};
 \item положительная спектральная щель: \textbf{получена};
 \item устойчивость по положительным весам: \textbf{48 из 48};
 \item незавешенный представитель: \textbf{существует};
 \item единственная метрика относительных скоростей: \textbf{не выведена};
 \item детальный баланс по следу выбирает веса: \textbf{нет};
 \item абсолютная физическая скорость: \textbf{не выведена};
 \item следующий тест: \textbf{нетривиальный KMS-селектор скоростей}.
\end{enumerate}
\end{protocol}

Машинный сертификат сохранён в
\path{s2t_v8_full_primitive_markov_generator_assembly_gate_results.json}.

\begin{thebibliography}{9}
\bibitem{v8-full-primitive-fagnola}
F.~Fagnola, C.~Mora,
\emph{On the relationship between a quantum Markov semigroup and its
representation via linear stochastic Schr\"odinger equations},
Open Systems \& Information Dynamics \textbf{21} (2014), 1450001;
arXiv:1405.6374.
\bibitem{v8-full-primitive-gao}
L.~Gao, C.~Rouz\'e,
\emph{Complete entropic inequalities for quantum Markov chains},
Archive for Rational Mechanics and Analysis \textbf{245} (2022),
183--238; arXiv:2102.04146.
\end{thebibliography}